\chapter*{Prerequisite bridge}
\addcontentsline{toc}{chapter}{Prerequisite bridge}

This short bridge is for readers who have not studied computer science. You need only ordinary arithmetic, careful reading, and the willingness to test claims.

\section*{Numbers and units}
Computers represent quantities using finite strings of bits. A bit has two possible states, conventionally written 0 and 1. Eight bits form a byte. A kilobyte is used in two different ways: decimal SI usage means 1,000 bytes, while the binary unit kibibyte means 1,024 bytes. Technical writing should state which convention it uses.

\section*{Functions as input-output rules}
In mathematics, a function assigns exactly one output to each allowed input. If $f(x)=2x+1$, then $f(3)=7$. In programming, a function may also change state, perform I/O, raise an exception, or return different results because of time or randomness. The mathematical analogy is useful, but it is not a complete definition of a program function.

\section*{Evidence and uncertainty}
A measurement is an observation produced by a procedure. It can be noisy, biased, incomplete, or affected by the act of measurement. A sample is not automatically representative of a population. A result can be compatible with a hypothesis without proving it. These distinctions reappear in usability studies, organisational evaluation, and machine learning.

\section*{The minimum command-line vocabulary}
The command line is a text interface for starting programs and inspecting files. A command has a current working directory, inputs, outputs, and an exit status. Version control records changes to files over time. The appendix introduces both tools carefully; the early chapters use commands only when they illuminate the program.

\plainlanguage{A technical system is easier to understand when you can state what goes in, what comes out, what changes in between, and what would count as failure. That is the habit this bridge is designed to establish.}

